% Chapter 3: Market Size - merged into ch02, this file kept minimal
% Projections table only
\chapter{Proyecciones de Demanda de Minerales a 5 Años}

\subsection{Proyección de Demanda de Cobre}

El cobre es el mineral más crítico para la transición energética. La producción primaria global podría alcanzar su pico alrededor de 2030 a $\sim$33 Mt antes de comenzar a declinar por agotamiento de yacimientos existentes.

\begin{keyinsightbox}[Déficit Estructural de Cobre]
Se proyecta un déficit persistente de oferta de cobre a partir de 2028--2030. Cada vehículo eléctrico requiere $\sim$80 kg de cobre (vs.\ $\sim$20 kg en vehículos convencionales). La demanda de cobre en infraestructura de carga de EVs crecerá >1.000\% para 2030 vs.\ niveles de 2020. El pipeline de nuevos proyectos es ``delgado'' — desarrollar un proyecto desde descubrimiento a producción toma 15--25 años.
\end{keyinsightbox}

\begin{table}[htbp]
\centering
\caption{Proyecciones de Demanda de Minerales Clave por País (2026--2031)}
\begin{tabular}{@{}llrrrr@{}}
\toprule
\textbf{Mineral} & \textbf{País Clave} & \textbf{2026} & \textbf{2028} & \textbf{2031} & \textbf{CAGR} \\
\midrule
Cobre & Chile & 5.8 Mt & 6.1 Mt & 6.5 Mt & 2.3\% \\
\rowcolor{tablealt} Cobre & Perú & 2.8 Mt & 3.1 Mt & 3.5 Mt & 4.6\% \\
Cobre & USA & 1.2 Mt & 1.4 Mt & 1.8 Mt & 8.4\% \\
\rowcolor{tablealt} Hierro & Brasil & 380 Mt & 395 Mt & 420 Mt & 2.0\% \\
Litio & Chile & 280 kt LCE & 350 kt & 450 kt & 9.9\% \\
\rowcolor{tablealt} Litio & Brasil & 45 kt LCE & 80 kt & 140 kt & 25.5\% \\
Litio & USA & 35 kt LCE & 60 kt & 100 kt & 23.4\% \\
\rowcolor{tablealt} Oro & Rusia & 340 t & 370 t & 410 t & 3.8\% \\
Cobre & Rusia & 1.0 Mt & 1.1 Mt & 1.2 Mt & 3.7\% \\
\rowcolor{tablealt} Tierras Raras & USA & 45 kt & 65 kt & 95 kt & 16.1\% \\
\bottomrule
\end{tabular}
\figuresource{S\&P Global \cite{spglobal2026}, Wood Mackenzie \cite{woodmac2026}, COCHILCO \cite{cochilco2026}, USGS, análisis propio}
\label{tab:mineral_demand}
\end{table}

\subsection{Escenarios de Precio del Cobre}

\begin{table}[htbp]
\centering
\caption{Escenarios de Precio del Cobre (USD/lb, 2026--2031)}
\begin{tabular}{@{}lrrrrr@{}}
\toprule
\textbf{Escenario} & \textbf{2026} & \textbf{2027} & \textbf{2028} & \textbf{2030} & \textbf{2031} \\
\midrule
Conservador & \$4.20 & \$4.10 & \$4.30 & \$4.50 & \$4.60 \\
\rowcolor{tablealt} Caso Base & \$4.50 & \$4.70 & \$5.00 & \$5.50 & \$5.80 \\
Optimista & \$4.80 & \$5.20 & \$5.80 & \$6.50 & \$7.00 \\
\bottomrule
\end{tabular}
\figuresource{Wood Mackenzie \cite{woodmac2026}, Bloomberg consensus, análisis propio}
\label{tab:copper_scenarios}
\end{table}

\subsection{Implicaciones para Consultoría de Ingeniería}

La demanda de servicios de consultoría minera está directamente correlacionada con el CAPEX minero. Con un pipeline combinado de >USD 215B en los cinco países objetivo:

\begin{itemize}
    \item \textbf{Estudios de factibilidad:} Representan típicamente 1--3\% del CAPEX total del proyecto.
    \item \textbf{Ingeniería EPCM:} 5--8\% del CAPEX total.
    \item \textbf{Optimización operacional:} Mercado recurrente, creciente con complejidad de operaciones.
    \item \textbf{Consultoría ESG/cierre:} Segmento de más rápido crecimiento (CAGR >10\%).
\end{itemize}

\begin{opportunitybox}[Oportunidad Estimada de Consultoría]
Pipeline total de CAPEX en países objetivo: $\sim$USD 215B. \\
Ratio consultoría/CAPEX típico: 3--6\% del valor del proyecto. \\
\textbf{Mercado direccionable estimado: USD 6.5--13B} en servicios de consultoría e ingeniería durante el período 2026--2034, solo en los cinco países objetivo.
\end{opportunitybox}
